\section{Conclusion}

\headingindent
During this internship, we investigated several ways of using model
distillation for OOD detection: classical output-level distillation,
distillation under pixel- and embedding-level perturbations, ensembles of
students trained on fixed feature subspaces, and feature denoising. The last
direction departed from classical distillation by asking a student to recover
a clean teacher representation from a corrupted one. Together, these
experiments examined whether the information lost, distorted, or disagreed
upon by a restricted student could provide a reliable indication that an
input does not belong to the training distribution.

The experiments show that such signals do exist. Teacher--student mismatch,
student confidence, ensemble disagreement, and feature reconstruction all
separated ID and OOD samples to some extent. Feature denoising was particularly
effective on several Far-OOD datasets, where the distribution shift is large.
Output-level distillation was generally more competitive on Near-OOD datasets,
where the student could still exploit class-related information. These results
suggest that the two families capture different aspects of distribution shift:
feature predictability is useful for large structural changes, whereas logits
and probabilities remain more sensitive to changes close to the teacher's
decision boundary. This qualitative Near/Far split is consistent with the
findings of Activation Subspaces \cite{zongur2025actsub}. That work reports
that classifier-insignificant activation components are especially useful for
Far-OOD detection, while decisive, class-related components are more useful
for Near-OOD detection.
% Although the methods are different, both sets of
% experiments indicate that weakly class-dependent representation information
% helps with large distribution shifts, whereas class-decisive information
% remains important when OOD samples are close to the ID distribution.

However, no student-based method consistently surpassed the teacher MSP
baseline across CIFAR-10, CIFAR-100, and ImageNet-200. The main difficulty was
Near-OOD detection, especially when ID and OOD samples shared similar visual
and semantic structure. The tested distillation objectives did not preserve a
sufficiently strong signal for these subtle differences, and their results
also remained below several established OOD methods.
% Consequently, the
% experiments do not validate the original idea in its current form: model
% distillation alone did not produce a detector that was both stronger than the
% teacher and stable across datasets.

In the final discussion with my supervisor, we considered whether the results
were promising enough to continue along the same path. We concluded that the
current formulation was not working sufficiently well to justify only scaling
up the same experiments and that this is a normal point in an exploratory research
project: once the main hypothesis has been tested and its limitations become
clear, the next step is usually to revise the hypothesis, modify the learning
objective, or pivot toward the part of the approach that showed the strongest
signal. Here, that could involve a more targeted study of Near-OOD failure
cases, a different objective for preserving subtle representation-level
differences, or a reformulation of feature denoising.

An additional and unexpected outcome of the internship concerned evaluation
methodology. The positive-class convention used for FPR@95--ID positive or
OOD positive--is still not standardized, despite the effort represented by
the OpenOOD benchmark. Even papers published at the same recent conference can
use different conventions. As demonstrated empirically in the OOD detection field,
changing the convention can produce a substantially lower FPR@95 from exactly
the same predictions, while AUROC remains unchanged. Without an explicit
definition of the positive class and threshold population, this difference
can be mistaken for a genuine improvement over another method.

% The internship therefore produced both a negative scientific result and a
% useful methodological one. The proposed distillation approach requires a
% substantial revision before it can become competitive, but the experiments
% identified its main failure mode and clarified which components retain useful
% OOD information. At the same time, the fixed evaluation pipeline, explicit
% metric conventions, and saved per-sample results provide a reproducible basis
% for any future reformulation of the idea.
